\section{Distant deletion comparisons}\label{sec:long-palm}

We now compare the deletion measures for distant cells of side $R$.
We prove \eqref{eq:distant-comparison} for the ordered $R$-cell pairs
$\mathcal P_R$ defined at the beginning of Section~\ref{sec:paths}.
The local construction of Section~\ref{sec:local-palm} cannot simply be
iterated: its route would have length growing with $K$, whereas even a
fixed loss at each step would destroy the estimate.  We instead retain
the original paths near both ends of a long random route.  In a product
of two likelihoods, the resulting errors occur only when the changed
portions of the two routes meet.  A geometric input makes these meetings
rare and controls their number exponentially.

\subsection{The geometric input and its application}

We use the following complete statement of \cite[Proposition~7.1]{D}.
It concerns a random set of lattice sites only.  The proof in
\cite[Section~7]{D} uses folded paths and the unpredictable-path
method of \cite{BPP}.

\begin{proposition}[Routes with rare encounters away from the common end]
\label{prop:route-input}
On $\mathcal L_K=(\Z/K\Z)^3$, let $d_\infty$ be the periodic maximum
distance.  A random datum $\mathcal D$ determines a bad set
$\mathcal B\subset\mathcal L_K$ satisfying
\[
 \Pp(A\subset\mathcal B)\le p^{|A|}\qquad(A\subset\mathcal L_K).
\]
Fix an integer contact radius $r_c\ge1$, an integer $q\ge0$, and $c_Q>0$.
Let $Q$ be a projected lattice cube of side at most $c_QR$, and let
$z$ be a site such that every $y\in Q$ has forward first-coordinate
displacement from $z$ in $[.29K,.46K]$.  There are constants
$p_0,\kappa,c,C,C_0>0$ and $R_0$, depending only on the fixed
parameters, with the following properties when $R\ge R_0$,
$K\ge C_0R$, and $p\le p_0$.

Put $H=\lfloor R^{1/2}\rfloor$.  There is an event $\mathcal E$,
determined by $\mathcal D,z,Q$, on which one can sample simple good
nearest-neighbor routes $\pi_y$ from $z$ to every $y\in Q$.
Conditional on $\mathcal D$, different samples use independent fresh
randomness.  The event $\mathcal E$ includes the requirement that the
$10R$-neighborhoods of $z$ and $Q$ contain no bad sites.

For $\pi=(v_0,\ldots,v_m)$ starting at $z$, put
\[
 a(\pi)=\max\{i:d_\infty(v_i,z)\le H\},\qquad
 \pi^{\rm ch}=\{v_i:i>a(\pi)\}.
\]
The initial portion $(v_0,\ldots,v_{a(\pi)})$ has at most $CH$
vertices.  The analogous portion read backwards from the terminal
site also has at most $CH$ vertices, and these two portions are
disjoint.  Let $\pi^{\rm ch,q}$ contain the vertices whose path index
is within $q$ of an index in $\pi^{\rm ch}$.  For independent samples
$\pi_y,\pi_{y'}$, set
\[
 J_c=\sum_{v\in\pi_y}\sum_{v'\in\pi_{y'}}
       \1_{\{d_\infty(v,v')\le r_c\}},\qquad
 J_*=\sum_{v\in\pi_y^{\rm ch,q}}
       \sum_{v'\in\pi_{y'}^{\rm ch,q}}
       \1_{\{d_\infty(v,v')\le r_c\}}.
\]
Uniformly in the specified endpoints,
\begin{equation}\label{eq:route-exponential}
 \E_{\mathcal D}\bigl[\1_{\mathcal E}
       \E_{\rm routes}(e^{\kappa J_c}\mid\mathcal D)\bigr]\le C.
\end{equation}
For every deterministic probability mass function $\omega$ on
$Q\times Q$ with $\omega(y,y')\le C_\omega R^{-6}$,
\begin{equation}\label{eq:route-rare}
 \sum_{y,y'\in Q}\omega(y,y')
   \E_{\mathcal D}\bigl[\1_{\mathcal E}
       \Pp_{\rm routes}(J_*>0\mid\mathcal D)\bigr]
 \le CC_\omega R^{-c_g},\qquad
 c_g=\frac12\left(\frac{\log2}{\log3}-\frac12\right)>0.
\end{equation}
Both estimates hold with a common terminal endpoint, reading the same
route construction in reverse order.  Finally, if $\mathcal M$ is
the stated bad-free margin event, then
\begin{equation}\label{eq:route-rejection}
 \Pp(\mathcal M\setminus\mathcal E)
 \le K^3(Ap)^{\lfloor K/(10(r_c+2))\rfloor}+CR^3e^{-cR},
\end{equation}
where $A$ is fixed and $Ap_0<1$.
\end{proposition}

We check the hypotheses and fix how this result is used.  The bad sites
are those of Section~\ref{sec:paths}.  Equation~\eqref{eq:good-prob}
gives the required joint bound with any sufficiently small fixed $p$,
by increasing the threshold for $D$.  All possible mean-sites of tags
whose midpoint lies in $v$ belong to a fixed enlargement $Q_v$ of that
cell, of side $O_x(R)$; this follows from the individual tubes.  The
same holds for $w'$.  Since $K/R\to\infty$ in the thermodynamic limit,
the displacements from any site of $Q_v$ to any site of $Q_{w'}$
belong to the wider interval in Proposition~\ref{prop:route-input}.

Choose $r_c$ large enough, in terms of $b_0$, to contain every possible
cross-route interaction between paths in their individual tubes, as
well as the bounded offsets between sites and means.  Take $q=2$ to
allow the adjacent labels used by a paired insertion.  These choices
are fixed once $x$ is fixed.  We retain a bad-free margin of radius
$C_xR$ around both cells and impose the route-availability events for
every endpoint site in their enlarged cubes, for both orientations.
There are only polynomially many such sites in $R$.  Thus
\eqref{eq:route-rejection} bounds the additional rejection by
$o_K(1)+C_xR^{C_x}e^{-c_xR}$.  This extra retention can only decrease
the nonnegative quantities in
\eqref{eq:route-exponential}--\eqref{eq:route-rare}.

We always sample the routes by the construction of
Proposition~\ref{prop:route-input}, conditional on the datum and the
endpoint sites.  In a likelihood product, the datum is shared and the
route randomizations are independent.  In particular, we do not
condition their law on a subsequent successful path insertion.
The exponent $c_g$ is independent of $x$, although the constants and
$\kappa$ may depend on $x$.

\subsection{A common measure with retained end portions}

Use the tag weight \eqref{eq:tag-pair-weight} for $(v,w')\in\mathcal P_R$,
retained data, and distinct tags $i_0,i_m\in I$.  On the sampled route
from their mean-sites choose one interior label at each intermediate
site, uniformly and independently given the tags and sites, excluding
the two tags.  The resulting list $i_0,\ldots,i_m$ has distinct labels;
each interior choice has atom at most $C_x/(Df_0)$.  Successive means
are at bounded absolute distance.

Write the initial silent portion of the route as $k<a_0$, and its
terminal silent portion as $k>a_1$.  Thus $a_0$ is one plus the last
visit within distance $H$ of the initial site, while $a_1$ is one less
than the first index in the terminal portion.  Both tags lie in these
portions, and the portions are disjoint.  Keep all unlisted paths in
both copies.  Keep the prefix originals $k<a_0$ on the left, and the
suffix originals $k>a_1$ on the right.  The remaining paths are
generated as follows, pairing left label $i_k$ with right label
$i_{k-1}$ for $1\le k\le m$.

\begin{enumerate}
\item For $a_0\le k\le a_1+1$, propose a common uniform midpoint
in a unit cell within bounded distance of both means and propose the
two paths independently conditional on that midpoint.  Condition
jointly on their individual cutoffs, the checks
\eqref{eq:soft-check} against retained and previously generated paths
in their respective copies, and midpoint separation by $s_1$ from
those opponents.  These middle pairs are probability kernels:
Proposition~\ref{prop:insertion}, Assertion~2, bounds each success
probability below by $1/2$.  Use any fixed order for them.

\item In every other pair, one path is a retained original.  Pin the
other path to its midpoint $y$.  Give this step mass zero unless $y$
is within $8b_0$ of the new label's mean and is $s_1$-separated from
all opponent midpoints already present in the copy receiving the
new path.  If these conditions hold, propose the bridge pinned at
$y$, impose its individual cutoffs and checks
\eqref{eq:soft-check}, and divide by
\begin{equation}\label{eq:pin-denominator}
 \max\{1/2,p_{\rm pin}\},
\end{equation}
where $p_{\rm pin}$ is their success probability under the pinned
proposal.  This is a subprobability kernel.  Perform these steps
after all middle pairs, in fixed orders.  Prefix pins are inserted
on the right and suffix pins on the left; neither family uses paths
generated by the other family.
\end{enumerate}

The final copies satisfy all conditional reference checks.  Every
left midpoint except that of $i_0$ agrees with its paired right
midpoint, and every right midpoint except that of $i_m$ is paired.
Deleting these respective tags and using the common ordering of the
bath coordinates therefore gives one measure on the bath with lifts
to both augmented references.  Figure~\ref{fig:silent-transport}
shows the pairing and the retained end portions.

\begin{figure}[tb]
\centering
\begin{tikzpicture}[x=.92cm,y=1.45cm,>=Latex,
 particle/.style={circle,draw,minimum size=7mm,inner sep=1pt,font=\small},
 made/.style={particle,fill=blue!12,draw=blue!65!black},
 pinned/.style={particle,fill=orange!18,draw=orange!70!black},
 old/.style={particle,fill=gray!15,draw=gray!65!black}]
\node[anchor=east] at (-.6,1) {left};
\node[anchor=east] at (-.6,0) {right};
\foreach \k in {0,1}{\node[old] (L\k) at (\k,1) {$i_\k$};}
\foreach \k in {2,3,4,5}{\node[made] (L\k) at (\k,1) {$i_\k$};}
\node[pinned] (L6) at (6,1) {$i_6$};
\node[pinned] (R0) at (0,0) {$i_0$};
\foreach \k in {1,2,3,4}{\node[made] (R\k) at (\k,0) {$i_\k$};}
\foreach \k in {5,6}{\node[old] (R\k) at (\k,0) {$i_\k$};}
\foreach \k/\j in {2/1,3/2,4/3,5/4}{\draw[blue!65!black,thick] (L\k)--(R\j);}
\draw[orange!70!black,thick,->] (L1)--(R0);
\draw[orange!70!black,thick,->] (R5)--(L6);
\draw[red!65!black,thick] (L0.south west)--(L0.north east);
\draw[red!65!black,thick] (R6.south west)--(R6.north east);
\draw[decorate,decoration={brace,amplitude=4pt},gray!65!black]
 (-.35,1.45)--(1.35,1.45) node[midway,above=5pt,font=\small] {silent prefix};
\draw[decorate,decoration={brace,amplitude=4pt},gray!65!black]
 (4.65,1.45)--(6.35,1.45) node[midway,above=5pt,font=\small] {silent suffix};
\node[anchor=north,font=\small,align=center] at (3,-.55)
 {gray: retained originals\qquad blue: jointly generated pairs\\
 orange: paths pinned to a retained midpoint};
\end{tikzpicture}
\caption{A schematic long comparison with $a_0=2$ and $a_1=4$.
The left prefix and right suffix retain their original paths. Central
pairs receive common proposed midpoints; an end pin copies the midpoint
of its retained partner in the other copy. The arrows indicate the
source of a pin, not equality of full paths. The changed left set is
$\{i_2,\ldots,i_6\}$, which excludes the common first tag. Deleting the
crossed tags again leaves identical midpoint multisets.}
\label{fig:silent-transport}
\end{figure}


The silent portions have two purposes.  They keep the common tag
outside the changed set in a likelihood product, and they ensure that
the changed sets start far enough from their common end for
\eqref{eq:route-rare} to apply.  The opposite-copy pins preserve the
matching of bath midpoints across these retained portions.

\begin{lemma}\label{lem:long-mass}
The common measure just constructed has mass satisfying
\begin{equation}\label{eq:long-mass}
 \operatorname{mass}\ge \E[A_vA_{w'}]-CD^{-1-c'}-o_K(1)
\end{equation}
for some $c'>0$ and sufficiently small fixed $w,x$.
\end{lemma}
\begin{proof}
The tag omissions, the bad-free margins, and the route rejections
have the stated cost by \eqref{eq:good-prob},
\eqref{eq:tag-omission}, and \eqref{eq:route-rejection}; the margins
contain $O_x(R^3)$ sites, and extra relative count factors cost only
fixed powers $D^x$, apart from negligible count tails.  Impose also
the full-path bounds of Section~\ref{sec:paths} on the original
sample in the endpoint neighborhoods of radius $O_x(R)$.  Their
failure has smaller cost, by the last line of
\eqref{eq:good-prob}.  This event is used only to estimate the mass;
we do not condition the reference bridges on it.

On this original-sample event, all pins lie in range: the retained
original has displacement at most $b_0$ from its endpoints, and
successive means are at bounded distance.  Whenever the required
midpoint separations hold, every $p_{\rm pin}\ge3/4$ by
Proposition~\ref{prop:insertion}, Assertion~3.  Indeed its opponents
are the original collection, with some originals deleted and some
generated paths added.  Route simplicity and the individual tubes
bound the number of added paths relevant to any fixed neighborhood
independently of the route length.  Consequently all such pin
kernels have mass one.

Separation from a middle path was already imposed when that middle
pair was generated: in the other copy its common midpoint was
tested against the retained original which supplies the pin.  The
two silent ends are far apart.  Every remaining separation failure
therefore entails two original midpoints at distance at most $s_1$
in an endpoint neighborhood.  Equation~\eqref{eq:pair-bound} bounds
their expected ordered count in each fixed site neighborhood by
$C_xD^2s_1^3$.

The choice of an interior label has atom at most $C_x/(Df_0)$ even
after fixing the data and tags.  A route visits a site at most once,
so we can sum this bound over all endpoint neighborhoods, of total
volume $O_x(R^3)$, without conditioning any count on route
availability.  A tag appearing in such a test supplies at least the
same weight through its divisor.  The other tag factors are bounded
by relative counts.  More explicitly, restrict all unit counts in
these neighborhoods to $CD$; by
\eqref{eq:count-tails}--\eqref{eq:count-mgf}, the discarded parts
contribute less than every inverse power of $D$, even after fixed
count products and polynomially many site sums.  The remaining loss
is at most
\[
 C_xR^3D^{Cx}\frac{D^2s_1^3}{Df_0}.
\]
Since $D^2s_1^3=D^{-2/5}$ and $R\asymp D^w$, this is
$O(D^{-1-c'})$ for sufficiently small $w,x$.  Middle steps never
lose mass.  Summing the tag weights proves \eqref{eq:long-mass}.
\end{proof}

\subsection{Conditional likelihoods and removal of cross checks}

We next bound the second moment of a lift.  Unlike a uniform bound
on its density, this estimate will have no factor depending on route
length.  Treat the left copy as the observed, or primary, copy.  Its
changed labels form
\[
 S=\{i_k:k\ge a_0\}.
\]
Fix the datum, tags, route and list, and condition on the paths $o$
outside $S$.  With $B_S$ the product plain bridge law and $c_S$ the
conditional reference indicator, write
\begin{equation}\label{eq:long-kernel}
 M_S(o;\dd z)=\int c_S(y;o)h_{w'i_m}(y)
                    K_S(y,o;\dd z)\,B_S(\dd y),\qquad
 Z_S(o)=\int c_S(y;o)\,B_S(\dd y).
\end{equation}
Here $y$ denotes latent original paths, $z$ the primary outputs,
and $K_S$ the subprobability kernel of all generation steps.  The
likelihood, before the divisor $T_{w'}$, is
$L_S=\dd M_S/\dd B_S$ on admissible reference outputs: the input
and output conditional normalizer $Z_S$ cancels exactly.  Absolute
continuity is verified below, including for the pins.

For an upper bound we may integrate out the pins generated on the
right at the silent prefix and replace their total mass by one.
No primary middle or pin step uses those paths.  We use the same
notation $M_S$ for this enlarged measure when taking upper bounds.

Take a second, independently sampled route and list with the same
first tag $i_0=i$.  Denote its changed set by $T$ and the two terminal
tags by $j_1,j_2$.  Condition now on all paths outside $U=S\cup T$,
as well as the datum.  In the $S$ construction remove all checks
against outside labels in $T\setminus S$, both from the latent
reference and from every proposal and success normalization.
In the $T$ construction remove checks against $S\setminus T$.
Write $A_S^-,Z_S^-,M_S^-$ for the resulting reference indicator,
normalizer, and unnormalized output measure, and similarly for $T$.
These two constructions now use only their own latent and generated
variables and the fixed paths outside $U$.

This independence also holds at the edge of a silent prefix.  Its
right pins have been integrated out.  If an old prefix path is
removed as an outside obstacle, no other part of the construction
uses that trajectory.  A right middle partner uses the fixed
endpoints of its label, not its old trajectory.  Routes and label
choices use only data, sites, and tag identities, all fixed here.

Let $Z_U=\int c_U\,\dd B_U$ be the union reference normalizer.
List all checks affected by the removals: reference checks and
checks of latent or generated paths against removed old obstacles,
including those in the success probabilities.  Every such check
requires two route vertices within $r_c$, allowing the fixed step
offsets in $J_*$.  There are at most $C_xJ_*$ such checks, with
bounded multiplicity.  The simplicity of each route gives this
bound regardless of its length.  Only distinct-label checks are
listed; in particular, a proposal never gains a check against its
own old trajectory if that trajectory was not an obstacle originally.

Let $B_{ST}$ be the sum, with these bounded multiplicities, of the
corresponding costs $\beta_{lk}$ from \eqref{eq:beta}.  For pin checks
also sum their costs $m(y,\omega)$ from
\eqref{eq:midpoint-cost}, writing the result as $M_{ST}'$.
The first quantity depends on the fixed endpoints and labels; the
second can depend on latent or generated midpoints.  Uniformly,
\begin{equation}\label{eq:cross-check-cost}
 B_{ST}+M_{ST}'\le C_xD^{-.19}J_c.
\end{equation}

The following estimate is the aim of the conditional comparison.  For
disjoint changed sets it preserves the original tag product, with an
error charged only to cross-route checks.  For overlapping sets it
charges a density bound only to the shared output labels.  Put
$G_D=C_xD^{C\sqrt x}$, with the constants sufficiently large.

\begin{lemma}[Conditional product bound]\label{lem:long-product}
Conditional on $\mathcal D$ and the paths outside $U=S\cup T$,
the two likelihoods satisfy
\begin{align}
 S\cap T=\varnothing:\quad
 \E[L_SL_T\mid U^c]
 &\le \E[h_{w'j_1}h_{w'j_2}\mid U^c]
   +Ce^{C_xD^{-.19}J_c}
       \bigl(B_{ST}+C_xJ_*b_D(1+t^{-3/2})\bigr),
       \label{eq:long-product}\\
 |S\cap T|=q\ge1:\quad
 \E[L_SL_T\mid U^c]
 &\le e^{C_xD^{-.19}J_c}G_D^q.
       \label{eq:long-product-overlap}
\end{align}
\end{lemma}

We prove the lemma through three comparisons: first of the reference
normalizers, then of the generation kernels when cross checks are
removed, and finally of the primary output densities on shared labels.

\begin{lemma}\label{lem:normalizer-comparison}
For these conditional reference laws,
\begin{equation}\label{eq:long-normalizer}
 \frac{Z_S^-Z_T^-}{Z_U}
       \le C^{|S\cap T|}e^{CB_{ST}}.
\end{equation}
If $S\cap T=\varnothing$, then for any measurable $0\le F\le1$
of the latent originals,
\begin{equation}\label{eq:reference-difference}
 0\le \frac1{Z_U}\int F(A_S^-A_T^--c_U)\,\dd B_U
          \le CB_{ST}e^{CB_{ST}}.
\end{equation}
\end{lemma}
\begin{proof}
For disjoint sets, $Z_S^-Z_T^-$ is the reference integral with
the cross checks removed.  Add these checks one at a time.
Conditional single-label insertion gives density domination by a
constant times the corresponding plain bridge, uniformly with any
subset of checks retained.  A newly added check therefore fails
with probability at most $C\beta_{lk}$.  Since the maximum of these
costs tends to zero, the loss of normalizer is at most
$e^{C\beta_{lk}}$ at that step.  Multiplication proves
\eqref{eq:long-normalizer} in the disjoint case.  Summing the
integrals lost when the checks are restored gives
\eqref{eq:reference-difference}.  All normalizers are positive by
successive insertion; no lower bound uniform in the number of
labels is being used.

If $q=|S\cap T|>0$, first delete the shared labels from both
separate integrals, which increases them because the integrands
are indicators.  Compare the remaining disjoint integrals as
above.  Insert the $q$ shared labels into the union reference,
one at a time; each insertion costs at most a fixed inverse
success probability.  Their total expense is $C^q$, proving
\eqref{eq:long-normalizer}.
\end{proof}

There is also a comparison before any generated variables are
integrated.  The original two constructions are dominated by the
separate constructions times a factor bounded in both forms
\begin{equation}\label{eq:protocol-comparison}
 e^{C_xD^{-.19}J_c},\qquad
 1+Ce^{C_xD^{-.19}J_c}(B_{ST}+M_{ST}').
\end{equation}
To prove this, fix a history at an affected transition.  Removing
checks increases its success probability from $p$ to $p^-$.  For a
middle step $p\ge1/2$; for a pin the denominator is the function
$\max(1/2,p)$, which is Lipschitz and bounded below by $1/2$.
Thus the expense in the likelihood density is at most
\[
 \frac{\max(1/2,p^-)}{\max(1/2,p)}\le1+2(p^--p),
\]
with the same estimate for the uncapped middle denominators.  Removing the preliminary pin-separation indicator
only enlarges the integrand.  For a middle step, integrate the
removed checks against its common uniform midpoint proposal;
Proposition~\ref{prop:insertion} bounds the difference by the sum
of $C\beta_{lk}$, including failures of midpoint separation.  For
a pin step the same bound includes its costs
\eqref{eq:midpoint-cost}; the needed midpoint separation holds
on the support of the old integrand.  Products of these bounds,
and \eqref{eq:cross-check-cost}, prove
\eqref{eq:protocol-comparison}.  Indicators of fewer reference
checks only enlarge an integrand.  Every unaffected transition
has ratio exactly one.

\subsection{Primary output densities and the conditional product bound}

We need one density estimate when two changed sets share labels.
For any $q$ specified primary labels, the marginal of $M_S^-/Z_S^-$
has density at most
\begin{equation}\label{eq:primary-density}
 G_D^q\quad\hbox{against their product plain bridges}.
\end{equation}
To prove it, omit the terminal tag indicator, so the latent originals
have their normalized separate reference law.  Fix those originals,
denoted by $y$, and let $K_S^-$ denote the separate generation kernel.
Divide the specified primary indices into middle indices $A_{\rm mid}$
and pin indices $A_{\rm pin}$.  A primary pin with index $k$ uses the
midpoint of the retained original label $i_{k-1}$ in the right suffix.
Write $\mathcal R_k=\{\xi:|\xi-m_{i_k}|\le8b_0\}$ for its pin range
in the compatible lift.

Conditional on the generation history, a specified middle output
is bounded by $C_xD^{Cx}B_{i_k}$, by
\eqref{eq:local-proposal}.  A specified pin is bounded by
$2\1_{\mathcal R_k}(y_{i_{k-1}}(0))
B_{i_k}^{y_{i_{k-1}}(0)}$, where the superscript denotes a prescribed
midpoint.  These majorants no longer depend on earlier generated
paths.  Integrating the generation steps in reverse order against
a nonnegative test of the specified outputs therefore gives the
conditional marginal bound
\begin{equation}\label{eq:primary-conditional-majorant}
\begin{split}
 K_S^-(y,o;\dd z_A)\le{}&
 (C_xD^{Cx})^{|A_{\rm mid}|}2^{|A_{\rm pin}|}
 \prod_{k\in A_{\rm mid}}B_{i_k}(\dd z_k)\\
 &\times\prod_{k\in A_{\rm pin}}
 \1_{\mathcal R_k}(y_{i_{k-1}}(0))
 B_{i_k}^{y_{i_{k-1}}(0)}(\dd z_k).
\end{split}
\end{equation}
Here $A=A_{\rm mid}\cup A_{\rm pin}$, and $z_k$ denotes the primary
output path of label $i_k$.  Every unspecified step contributes at
most its mass, which is at most one.  Thus the displayed product
bound does not require independence of the accepted outputs.

The sources $i_{k-1}$ of the specified pins are distinct latent labels
in $S$.  Corollary~\ref{cor:bridge-domination} bounds their joint
latent marginal by $C^{|A_{\rm pin}|}$ times their product plain
bridges.  Keep their confinement indicators when making this
comparison, so all source midpoints use the chosen local lifts.
Dropping the remaining source cutoffs in those lifts only enlarges
the integral.  For a source $l=i_{k-1}$, evaluate the resulting
midpoint mixture by the measure identity
\begin{equation}\label{eq:pin-mixture-density}
 \int_{\mathcal R_k}g_l(\xi)B_{i_k}^{\xi}(\dd z_k)\,\dd\xi
 =\1_{\mathcal R_k}(z_k(0))
   \frac{g_l(z_k(0))}{g_{i_k}(z_k(0))}
       B_{i_k}(\dd z_k).
\end{equation}
The two means are at bounded distance, and the indicator restricts
the ratio to the primary pin range.  Thus
\eqref{eq:midpoint-ratio} bounds this ratio by
$C_xD^{C\sqrt x}$.  Applying this to each pin in
\eqref{eq:primary-conditional-majorant} proves
\eqref{eq:primary-density}.  The same argument with the original
checks in place proves absolute continuity of all primary outputs
in \eqref{eq:long-kernel}.  Only the primary marginal acquires this
density; a pair with a common midpoint is still singular relative
to two independent full bridges.

\begin{proof}[Proof of Lemma~\ref{lem:long-product}]
Integrate the enlarged output measures against the common union
reference, with divisor $Z_U$, and use
\eqref{eq:protocol-comparison}.  Dropping the union output
admissibility indicator increases the resulting nonnegative
integrals.

First suppose $S,T$ are disjoint.  The leading $1$ in the second
bound of \eqref{eq:protocol-comparison} contributes at most
\[
 \frac1{Z_U}\int h_{w'j_1}h_{w'j_2}A_S^-A_T^-\,\dd B_U:
\]
the two separate kernels use different outputs and have mass at
most one.  Equation~\eqref{eq:reference-difference} compares this
to the leading conditional expectation in
\eqref{eq:long-product}.  The $B_{ST}$ error is bounded with
\eqref{eq:long-normalizer}.

For a term of $M_{ST}'$, the pin midpoint belongs to a latent
source $l$ in one separate construction and is tested against an
output path $z_k$ in the other.  After the cross checks have been removed,
the two constructions have product law at the fixed outside paths.
Write $\lambda_S^-$ for the normalized separate latent reference
and $\nu_T^-=M_T^-/Z_T^-$ for the opposing output measure, of mass
at most one.  The source-side kernels also have mass at most one.
Dropping the source-side tag indicator and integrating those kernels bounds
the normalized cost by
\[
 \int \nu_T^-(\dd z)\int\lambda_S^-(\dd y)
   \frac{b_D}{s_1\vee d_K(y_l(0),z_k(0))}
 \le C\sup_{\xi}\int
   \frac{b_Dg_l(\eta)}{s_1\vee d_K(\eta,\xi)}\,\dd\eta.
\]
Here single-label insertion bounded the source midpoint marginal
by $Cg_l$ in its local lift.  For every fixed opposing midpoint
$\xi$,
\begin{equation}\label{eq:pin-cost-average}
 \int \frac{b_D}{s_1\vee d_K(\eta,\xi)}g_l(\eta)\,\dd\eta
          \le Cb_D(1+t^{-3/2}).
\end{equation}
Indeed $g_l$, or its periodization, is bounded by $Ct^{-3/2}$,
the inverse-distance function is integrable on a unit ball in
three dimensions, and outside that ball it is at most one.
There is no need to bound the density of the opposing midpoint.
There are at most $C_xJ_*$ terms.  The normalizer ratio in
\eqref{eq:long-normalizer} completes the disjoint estimate.

When $q=|S\cap T|>0$, use the exponential bound in
\eqref{eq:protocol-comparison}.  After division by $Z_S^-Z_T^-$,
the remaining integral is the product of the two normalized
separate output densities, identifying the common output
coordinates only on $S\cap T$.  Integrating the unshared outputs
leaves the product of the two marginals on these $q$ labels.
Bound one by \eqref{eq:primary-density}; the other has mass at
most one.  Equation~\eqref{eq:long-normalizer} proves
\eqref{eq:long-product-overlap}.

For clarity, these comparisons take place first between measures
on latent originals and generated proposals, with all success
indicators and denominators present.  Their domination passes
to the primary output densities.  In the disjoint estimate we
retain the latent pin source until applying
\eqref{eq:pin-cost-average}; disintegration over the absolutely
continuous primary outputs justifies this operation.  In the
overlap estimate the two latent integrations remain separate,
and density versions agree almost everywhere under the union
product bridge measure by Fubini's theorem.  At no point do we
assign a product-bridge density to a midpoint-matched pair.
\end{proof}

\subsection{Averaging the labels and route encounters}

Let $\ell_v$ be the left lift likelihood, with zero value on
omitted data and tags.  On the retained part it is the sum of
$\overline{L_S}/T_{w'}$ over second tags, where the bar averages
the route and interior label kernels.  Squaring it gives two
independent choices with common first tag $i$.  The augmented
reference supplies its factor $h_{vi}/T_v$ only once; this path
is fixed outside $U$.  The disjoint leading term in
\eqref{eq:long-product}, summed over choices, is bounded by
\begin{equation}\label{eq:long-leading-term}
 \E[A_vA_{w'}^2].
\end{equation}
Here the indicators are the original reference indicators.
Removing good-data and distinctness restrictions only increases
their nonnegative sum, and the route and label kernels have mass
one.

We detail the remaining sums because an estimate uniform in route
length is essential.  At every possible tag mean-site, the number
of candidate tags is at most $C_xDf_0$, by the good-data flagged
counts and the bounded endpoint offsets.  Since
$T_v=T_{w'}=f_0DR^3$, its normalized counting measure is dominated
by a deterministic site weight $C_xR^{-3}$ times a label
subprobability measure whose atoms are at most $C_x/(Df_0)$.
The same atom bound holds for every interior slot, using its
actual conditional probability law.  Route sampling depends on
the endpoint sites and datum, not on tag identities.

For a disjoint error, enumerate all slot pairs across the two
routes that are within the contact radius, including the fixed
step offsets.  This overcounts the affected checks by a bounded
factor.  Every such check has an averageable label, and
\eqref{eq:label-average} gives an average cost at most $C_xb_D$
for its $\beta$.  In detail, interior choices are independent
across routes given the tags; a label in the opposite slot can
be held fixed.  When a tag must instead be summed, its point
weight gives the same bound.  A prefix label used by the pair
at the change edge is assigned to its adjacent prefix slot.
The common first tag is outside both changed sets.  Exclusions
such as omitting a label's own old trajectory remove terms;
we may discard disjointness restrictions in these positive,
capped-cost sums.  Thus the entire disjoint error after label
averaging is at most
\begin{equation}\label{eq:averaged-disjoint-error}
 C_xe^{C_xD^{-.19}J_c}J_*b_D(1+t^{-3/2}),
\end{equation}
with the deterministic site weights still to be summed.

For overlapping sets, identical labels across the two lists
form a matching of their slots, since each individual list has
distinct labels.  The common first tag cannot participate.  A match
requires a contact, so the number $M_*$ of eligible slot pairs is
at most $C_xJ_*$ and at most $C_xJ_c$.  Put $p_D=C_x/(Df_0)$.
Prescribing $m'$ disjoint matches has total label weight at most
$p_D^{m'}$.  To see the joint bound, first fix the two terminal tags.
For every prescribed edge involving an interior slot, integrate one
of its interior choices using the atom bound $p_D$.  These choices
are independent given the tags, and disjoint matching edges use
distinct choices.  If a terminal-tag pair is prescribed, sum that
match last using the subprobability point bound.  This proves the
bound also when different edges match each tag to an interior slot
of the other list.

Write $q=|S\cap T|$, and let $\E_{\rm labels}$ denote the sum with
the interior probability kernels and terminal-tag subprobability
weights, at the fixed sites and data.  Expand $G_D^q$ as the product
of $1+(G_D-1)\1_{\rm match}$ over eligible pairs.  To restrict to
$q\ge1$, designate one actual match first.  The preceding joint
bound, with incompatible prescriptions contributing zero, gives
\[
 \E_{\rm labels}[\1_{\{q\ge1\}}G_D^q]
 \le M_*p_DG_D\exp\bigl(M_*p_D(G_D-1)\bigr).
\]
Combining this with \eqref{eq:long-product-overlap} bounds the
overlap contribution by
\begin{equation}\label{eq:averaged-overlap-error}
 C_x\frac{G_D}{Df_0}J_*
   \exp\left(C_xD^{-.19}J_c+C_x\frac{G_D}{Df_0}J_c\right).
\end{equation}
Only matched slots incur a cost; unused route vertices incur none.

The sum of \eqref{eq:averaged-disjoint-error} and
\eqref{eq:averaged-overlap-error} has the form
\begin{equation}\label{eq:route-error}
 C_xD^{-1+\zeta(x)}J_*e^{\delta_DJ_c},\qquad
 \zeta(x)\longrightarrow0\ (x\downarrow0),\quad
 \delta_D\longrightarrow0\ (D\to\infty),
\end{equation}
where the second limit is for a sufficiently small fixed $x$.
The factors $b_D(1+t^{-3/2})$ have only polylogarithmic losses,
and $G_D/(Df_0)=C_xD^{-1+x+C\sqrt x}$, so these choices are
possible.

We can now use the geometry.  For each fixed common site, the
deterministic weights of the two noncommon sites are bounded
by $C_xR^{-6}$.  Enlarge their sum to a fixed cube of side
comparable with $R$ and compare with its uniform average.
Equations~\eqref{eq:route-exponential}--\eqref{eq:route-rare}
and Cauchy--Schwarz give
\begin{align*}
 &\E_{\rm sites}\E\bigl[
       \1_{\rm retained}J_*e^{\delta_DJ_c}\bigr]\\
 &\quad\le
  \Bigl(\E_{\rm sites}\E[
       \1_{\rm retained}\1_{\{J_*>0\}}]\Bigr)^{1/2}
  \Bigl(\E_{\rm sites}\E[
       \1_{\rm retained}J_c^2e^{2\delta_DJ_c}]\Bigr)^{1/2}
 \le C_xR^{-c_g/2}.
\end{align*}
For the last factor take $D$ large enough that $2\delta_D<\kappa/2$
and absorb $J_c^2$ into $e^{\kappa J_c/2}$.  The common-site
weights, each bounded by $C_xR^{-3}$ on $O_x(R^3)$ sites, have
bounded total mass.  All these bounds precede the data average;
the site weights are not probabilities conditioned on availability.

Choose $x$ so small that $\zeta(x)<wc_g/4$.  Since
$R\asymp D^w$, \eqref{eq:route-error} has total expectation
$O(D^{-1-c''})$ for some $c''>0$.  The permissible bad-set base
and $\kappa$ can then be fixed for this $x$, and finally $D$ can
be increased.  The absolute positive exponent $c_g$ makes this
order of choices possible.  We have proved
\begin{equation}\label{eq:long-second}
 \int\ell_v^2\,\dd P_{{\rm aug},v}
 \le \E[A_vA_{w'}^2]+O(D^{-1-c''})+o_K(1),
\end{equation}
in particular after averaging over $\mathcal P_R$.

The right lift of the same common measure satisfies
\eqref{eq:long-second} with $v,w'$ exchanged.  To verify this,
make the right copy primary and read the indices in reverse,
while preserving the actual generation order.  Integrate out
the unobserved pins near its common end; none of the remaining
kernels uses them.  All conditional comparisons above allow
that order.  For route averaging, use the common-terminal
conclusions of Proposition~\ref{prop:route-input} for the same
route law, rather than a different coupling.

\subsection{Completion of the distant comparison}

The reference mass is $\E A_v$ and the lift mass is the common
mass.  Subtract twice \eqref{eq:long-mass} from
\eqref{eq:long-second} and add the reference mass.  This gives,
after the indicated cell-pair average,
\begin{equation}\label{eq:long-likelihood-error}
 \int(\ell_v-1)^2\,\dd P_{{\rm aug},v}
 \le \E[A_v(A_{w'}-1)^2]+o(D^{-1})+o_K(1).
\end{equation}
Up to negligible count-tail terms, $A_v\le CD^x$.
Equation~\eqref{eq:flag-moments}, together with the uniform
marginal distribution of each cell in $\mathcal P_R$, therefore
bounds the average of the main term by
\[
 C\bigl(D^{-1-w/6+x}+D^{-1+2x}R^{-3}\bigr)=o(D^{-1})
\]
for sufficiently small $x$ relative to $w$.  The same holds for
the right lift.

Let $q$ be the density of the common projected measure on the
bath.  Jensen's inequality under projection preserves both
squared likelihood bounds.  Since $P_v$ has density $s_v^2$,
\[
 \int|s_v-\sqrt q|^2\,\dd Y'
 =\int s_v^2\left|1-\sqrt{q/s_v^2}\right|^2\,\dd Y'
 \le\int s_v^2|1-q/s_v^2|^2\,\dd Y'.
\]
Use absolute continuity for the zero-density convention.
Applying the same inequality with $w'$ and then the squared
triangle inequality proves \eqref{eq:distant-comparison}.
